\documentclass[12pt]{article}

\usepackage{amsfonts}
\usepackage{amsmath}

\begin{document}

pag 7

\hfill

\textbf{Chapter 1. The Orbits of One-Dimensional Maps.}

\hfill

\textbf{1.1 Iteration of Functions and Examples of Dynamical Systems.}

\hfill

Dynamical systems is the study of how things change over time. Examples include
the growth of populations, the change in the weather, radioactive decay, mixing of
liquids and gases such as the ocean currents, motion of the planets, the interest in a
bank account. Ideally we would like to study these with continuously varying time
(what are called flows), but in this text we will simplify matters by only considering
discrete changes in time. For example, we might model a population by measuring
it daily. Suppose that xn is the number of members of a population on day n, where
$x_0$ is the initial population. We look for a function $f :\mathbb{R} \to \mathbb{R}$, (where $\mathbb{R}$ = set of all
real numbers), for which

\[
x_1 = f(x_0), x_2 = f(x_1) \text{ and generally } x_n = f(x_{n-1}), n= 1, 2, ...
\]

This leads to the iteration of functions in the following way:

\hfill

\textbf{Definition 1.1.1} Given $x_0 \in \mathbb{R}$, the orbit of $x_0$ under f is the set

\[
O(x_0) = {x_0, f(x_0), f^2(x_0), ...},
\]

where $f^2(x_0) = f(f(x_0)), f^3(x_0) = f(f^2(x_0))$, and continuing indefinitely, so that

\[
f^n(x) = f \circ f \circ f \circ \cdots \circ f(x); (n-times composition).
\]

Set $x_n = f^n(x_0), x_1 = f(x_0), x_2 = f^2(x_0)$, so that in general

\[
x_{n+1} = f^{n+1}(x_0) = f(f^n(x_0)) = f(x_n).
\]


More generally, f may be defined on some subinterval I of $\mathbb{R}$, but in order for the
iterates of $x \in I$ under f to be defined, we need the range of f to be contained in I,
so $f: I \to I$.

\hfill

This is what we call the iteration of one-dimensional maps (as opposed to higher
dimensional maps $f: \mathbb{R}^n \to \mathbb{R}^n$, n > 1, which will be studied briefly in a later
chapter).

\hfill

\textbf{Definition 1.1.2} A (one dimensional) dynamical system is a function $f: I \to I$
where I is some subinterval of $\mathbb{R}$.

\hfill

Given such a function f, equations of the form $x_{n+1} = f(x_n)$ are examples of
difference equations. These arise in the types of examples we mentioned above. For example in biology xn may represent the number of bacteria in a culture after n hours.
There is an obvious correspondence between one-dimensional maps and these difference
equations. For example, a difference equation commonly used for calculating
square roots:

\[
x_{n+1} = \frac{1}{2}(x_n + \frac{2}{x_n}
\]
 
corresponds to the function $f(x) = \frac{1}{2} (x + \frac{2}{x_n})$. If we start by setting $x_0 = 2$ (or in
fact any real number), and then find $x_1, x_2$ etc., we get a sequence which rapidly
approaches $\sqrt{2}$ (see page 9 of Sternberg [63]). One of the issues we examine is what
exactly is happening here.

\hfill

\textbf{Examples of Dynamical Systems 1.1.3}

\hfill

\begin{enumerate}
\item \textbf{The Trigonometric Functions} Consider the iterations of the trigonometric
functions starting with $f: \mathbb{R} \to \mathbb{R}$, f(x) = sin(x). Select $x_0 \in \mathbb{R}$ at random, e.g.,
$x_0 = 2$ and set $x_{n+1} = sin(x_n)$, n = 0, 1, 2, ... Can you guess what happens to $x_n$ as
n increases? One way to investigate this type of dynamical system is to enter 2 into
our calculator, then repeatedly press the 2nd, Answer and Sin keys (you will need
to do this many times to get a good idea - it may be easier to use Mathematica, or
some similar computer algebra system). Do the same, but replace the sine function
with the cosine function. How do we explain what appears to be happening in each
case? These are questions that we aim to answer quite soon.

\item Linear Maps Probably the simplest dynamical system (and least interesting from
a chaotic dynamical point of view), for population growth arises from the iteration of
linear maps: maps of the form $f(x) = a \cdot x$. Suppose that $x_n$ = size of a population
at time n, with the property

\[
x_{n+1} = a \cdot x_n
\]

for some constant a > 0. This is an example of a linear model for the growth of the
population.

\hfill

If the initial population is $x_0 > 0$, then $x_1 = ax_0, x_2 = ax_1 = a^2x_0$, and in general
$x_n = a^nx_0$ for n = 0, 1, 2, . . .. This is the exact solution (or closed form solution)
to the difference equation $x_{n+1} = a \cdot x_n$. Clearly f(x) = ax is the corresponding
dynamical system. We can use the solution to determine the long term behavior of
the population:

$x_n$ is very well behaved since:

\begin{enumerate}
\item  if a > 1, then $x_n \to infty$ as $n \to \infty$,

\item  if 0 < a< 1 then $x_n \to 0$ as $n \to \infty$ (i.e., the population becomes extinct),
\item if a = 1, then the population remains unchanged.
\end{enumerate}

\hfill

\textbf{3. Affine maps} These are functions $f: \mathbb{R}
 \to \mathbb{R}$ of the form f(x) = ax + b ($a \neq 0$),
for constants a and b. Consider the iterates of such maps:

\[
f^2(x) = f(ax + b) = a(ax+ b) + b = a^2x + ab + b,
\]
\[
f^3(x) = a^3x + a^2b + ab + b,
\]
\[
f^4(x) = a^4x + a^3b + a^2b + ab + b,
\]

and generally

\[
f^n(x) = a^nx + a^{n-1}b + a^{n-2}b + \cdots + ab + b.
\]


Let $x_0 \in  \mathbb{R}$ and set $x_n = f^n(x_0)$, then we have

\[
\begin{matrix}
x_n & = & a^n x_0 + (a^{n-1} + a^{n-2} + \cdots + a + 1)b \\ 
& = & a^n x_0 + b \left ( \frac{a^n - 1}{a-1} \right ), \text{ if } a \neq 0 \\
\end{matrix}
\]

or

\[
x_n = \left ( x_0 + \frac{b}{a-1} \right ) a^n + \frac{b}{1-a}, \text{ if } a \neq 1
\]

is the closed form solution in this case (here we have used the formula for the sum of
a finite geometric series:

\[
\sum_{k=0}^{n-1} r^k = \frac{r^n -1}{r-1}
\]

when $r \neq 1$). If a = 1, the solution is $x_n = x_0 + nb$ .
We can use these equations to determine the long term behavior of xn. We see
that:

\begin{enumerate}
\item if $|a| < 1$ then $a^n \to 0$ as $n \to \infti$, so that

\[
lim_{n \to \infty} x_n = \frac{b}{1-a}
\]

\item  if $a > 1$, then $lim_{n \to \infty} x_n = \infty$ for b, $x_0 > 0$,

\item  if a = 1, then $lim_{n \to \infty = \infty$} if $b > 0$.
The limit won’t exist if a  −1.
\end{enumerate}

\item \textbf{The Logistic Map} $L{\mu}: \mathbb{R} \to \mathbb{R}, L_{\mu}(x) = \mux(1 − x)$ was introduced to model
a certain type of population growth (see [41]). Here $\mu$ is a real parameter which
is fixed. Note that if $0 < \mu \leq 4$ then $L_{\mu}$ is a dynamical system of the interval
[0, 1], i.e. $L{\mu}: [0, 1] \to [0, 1]$). For example, when $\mu = 4, L_4(x) = 4x(1 − x)$, with
$L_4([0, 1]) = [0, 1]$ with graph given below. If $\mu > 4$, then $L_{\mu}$ is no longer a dynamical
system of [0, 1] as $L_{\mu}([0, 1])$ is not a subset of [0, 1].

\hfill

\textbf{Recurrence Relations 1.1.4}

\hfill

Many sequences can be defined recursively by specifying the first term (or two),
and then stating a general rule which specifies how to obtain the nth term from the
(n − 1)th term (or other additional terms), and using mathematical induction to see
that the sequence is “well defined” for every $n \in  \mathbb{Z}^{+}$ = {1, 2, 3, ...}. For example,
n! = n-factorial can be defined in this way by specifying 0! = 1, and n! = n · (n −1)!,
for $n \in \mathbb{Z}^{+}$. The Fibonacci sequence $F_n$, can be defined by setting

\[
F_0 = 1, F_1 = 1 and F_{n+1} = F_n + F_{n−1}, for n \in \mathbb{Z}^{+}
\]

so that $F_2 = 2, F_3 = 5$ etc.

\hfill

The orbit of a point $x_0 \in \mathbb{R}$ under a function f is then defined recursively as follows:

we are given the starting value x0, and we set

\[
x_n = f(x_{n−1}), for n \in \mathbb{Z}^{+}.
\]

The principle of mathematical induction then tells that xn is defined for every n =  0,
since it is defined for n = 0, and assuming it has been defined for k = n − 1 then
$x_n = f(x_{n−1})$ defines it for k = n.

\hfill

Ideally, given a recursively defined sequence xn, we would like to have a specific
formula for xn in terms of elementary functions (so called closed form solution), but
this is often very difficult, or impossible to achieve. In the case of affine maps and
certain logistic maps, there is a closed form solution. One can use the former to study
problems of the following type:

\hfill

\textbf{Example 1.1.5} An amount \$ T is deposited in your bank account at the end of each
month. The interest is r% per period. Find the amount A(n) accumulated at the
end of n months (Assume A(0) = T).

\hfill

\textbf{Answer}. A(n) satisfies the difference equation

\[
A(n + 1) = A(n) + A(n)r + T, where A(0) = T,
\]

or

\[
A(n + 1) = A(n)(1 + r) + T,
\]

so we set $x_0 = T, a = 1+r$ and b = T in the formula of Example 1.1.3, then the
solution is

\[
A(n) = (1+r)^nT + T(\frac{(1 + r)n − 1}{1+r-1})
\]
\[
= (1+r)^nT + \frac{T}{r}((1+r)^n -1 )
\]

pag11

\hfill


\textbf{Remark 1.1.6} It is conjectured that closed form solutions for the difference equation
arising from the logistic map are only possible when $\mu = ± 2,4$ (see Exercises 1.1 #
3 for the cases where $\mu = 2, 4$, and # 7 for the case where $\mu = -2$ and also [67] for a
discussion of this conjecture).

\hfill

\textbf{Exercises 1.1}

\hfill

\begin{enumerate}
\item If $L_{\mu}(x) = \mux(1 − x)$ is the logistic map, calculate $L_{\mu}^2$ and $L^3_{\mu}(x)$.

\item Use Example 1.1.3 for affine maps to find the solutions to the difference equations:

\begin{enumerate}
\item $x_{n+1} - \frac{x_n}{3} = 2, x_0=2
\] 

\item $x_{n+1} + 3x_n = 4, x_0 = -1.$
\end{enumerate}

\item A logistic difference equation is one of the form $x_{n+1} = \mu x_n(1-x_n)$ for some fixed
$\mu \in \mathbb{R}$. Find exact solutions to the logistic equations:

\begin{enumerate}
\item $x_{n+1} = 2x_n(1 − x_n)$. Hint: use the substitution $x_n = (1 − y_n)/2$ to transform the
equation into a simpler equation that is easily solved.
\item $x_{n+1} = 4x_n(1 - x_n)$. Hint: set $x_n = sin^2(\theta_n)$ and simplify to get an equation that
is easily solved.
\end{enumerate}

\item You borrow \$ P at r % per annum, and pay off $M at the end of each subsequent
month. Write down a difference equation for the amount owing A(n) at the end of
each month (so A(0) = P). Solve the equation to find a closed form for A(n). If
P = 100, 000, M = 1000 and r = 4, after how long will the loan be paid off?

\item If 

\[
T_{\mu}(x) =  \left\{\begin{matrix}
\mu x & \text{if} & 0 \leq x < 1/2  \\
\mu(1-x) & \text{if} & 1/2 \leq x <1  \\
\end{matrix}\right.
\]

show that $T_{\mu}$ is a dynamical system of
[0, 1] for $\mu 2 (0, 2]$

\item Let $f: \mathbb{R}  \to \mathbb{R} $ be the function defined below. For each of the intervals given,
determine whether f can be considered as a dynamical system $f : I \to  I$:

\begin{enumerate}
\item $f(x) = x^3 − 3x$,

\hfill

(i) I = [−1, 1], (ii) I = [−2, 2].

\item $f(x) = 2x^3 − 6x,$

\hfill

(i) I = [−1, 1], (ii) $I = [- \sqrt{\frac{7}{2}}, \sqrt{\frac{7}{2}}],$
 (iii) I = [−4, 4].
 
\item For the following functions, find $f^2(x), f^3(x)$ and a general formula for $f^n(x)$:

\[
(i) f(x)=x^2, (ii) f(x) = |x+1|, (iii) f(x) = \left\{\begin{matrix}
2x & \text{ if }  & 0 \leq x < 1/2 \\
2x-1 & \text{ if }  & 1/2 \leq x < 1 \\
\end{matrix}\right.
\]

\item Use mathematical induction to show that if $f(x) = \frac{2}{x+1}$, then

\[
f^n(x) = \frac{2^n(x + 2) + (-1)^n(2x − 2)}{2^n(x + 2) - (-1)^n(x - 1)}
\]

\item Show that a closed form solution to the logistic difference equation when $\mu = -2$
is given by
\end{enumerate}

\{
xn =
1
2
 
1 − f
 
rnf−1(1 − 2x0)
  
, where r = −2 and f( ) = 2 cos
 
  −
p
3 
3
!
.
(Hint: Set xn =
1 − f( n)
2
and use steps similar to 3(ii)).

1.2 Newton’s Method and Fixed Points
Given a differentiable function f : R ! R, Newton’s method often allows us to
find good approximations to zeros of f(x), i.e., approximate solutions to the equation
f(x) = 0. The idea is to start with a first approximation x0 and look at the tangent
line to f(x) at the point (x0, f(x0)). Suppose this line intersects the x-axis at x1,
then we can check that
x1 = x0 −
f(x0)
f0(x0)
, if f0(x0) 6= 0.
For reasons that we will make clear, it is often the case that x1 is a better approximation
to the zero x = c than x0 was. In other words, Newton’s method is an algorithm
for finding approximations to zeros of a function f(x). The algorithm gives rise to a
difference equation
xn+1 = xn −
f(xn)
f0(xn)
,
where x0 is a first approximation to a zero of f(x). The corresponding real function
is
Nf (x) = x −
f(x)
f0(x)
, (the Newton function).
Note that if f(x) = x2 − a, then f0(x) = 2x and
Nf (x) = x −
x2 − a
2x
=
1
2
 
x +
a
x
 
,
so that
xn+1 =
1
2
 
xn +
a
xn
 
,
the difference equation we mentioned in Section 1.1 that is used for approximating p
2 when a = 2.

Note that we are looking for where f(x) = 0, and this happens if and only if
Nf (x) = x.
Definition 1.2.1 For a function f : R ! R, a point c 2 R for which f(c) = c is called
a fixed point of f. It is a point where the graph of f(x) intersects the line y = x. We
denote the set of fixed points of f by Fix(f).
Fixed points occur where the graph of f(x) intersects the line y = x.
Example 1.2.2 Suppose that f(x) = x2, then x2 = x gives x(x−1) = 0, so has fixed
points c = 0 and c = 1, so Fix(f) = {0, 1}. If f(x) = x3 − x, then x3 − x = x gives
x(x2 − 2) = 0, so the fixed points are c = 0 and c = ±
p
2, Fix(f) = {0,±
p
2}.

The logistic map f(x) = 4x(1 − x) = 4x − 4x2, 0   x   1 has the properties:
f(0) = 0, (a fixed point), f(1) = 0, a maximum at x = 1/2 (with f(1/2) = 1).
Solving f(x) = x gives 4x − 4x2 = x, so 4x2 = 3x, so the fixed points are c = 0 and
c = 3/4.
This map has what we call eventual fixed points: f(1) = 0 and f(0) = 0, so we say
that c = 1 is eventually fixed. Also f(1/4) = 3/4, so c = 1/4 is eventually fixed, as is
c = (2+
p
3)/4.
Definition 1.2.3 x  2 R is an eventual fixed point of f(x) if there exists a fixed point
c of f(x) and r 2 Z+ satisfying fr(x ) = c, but fs(x ) 6= c for 0 < s < r.
Example 1.2.4 The Tent Map
Define a function T : [0, 1] ! [0, 1] by
T(x) =
 
2x : 0   x   1/2
2(1 − x) :1/2 < x   1
= 1 − 2|x − 1/2|.
T(x) is called the tent map, it has fixed points: T(0) = 0 and T(2/3) = 2/3. Since
T(1/4) = 1/2, T(1/2) = 1 and T(1) = 0, 1/4, 1/2 and 1 are eventually fixed. It is not
difficult to see that there are many other eventually fixed points (infinitely many).
0.0 0.2 0.4 0.6 0.8 1.0
0.2
0.4
0.6
0.8
1.0
Note that some maps do not have fixed points:
Example 1.2.5 f(x) = x2 + 1 has no fixed points.

f(x) = x2 + 1 does not intersect the line y = x.
Question 1.2.6 For what values of a 2 R does Qa(x) = x2 + a have fixed points?
We see below that certain functions always have fixed points:
Theorem 1.2.7 Let f : I ! I be a continuous function, where I = [a, b], a < b is a
bounded interval. Then f(x) has a fixed point c 2 I.
Proof. Set g(x) = f(x) − x. We may assume that f(a) 6= a and f(b) 6= b, so that
f(a) > a and f(b) < b.
The graph of f(x) always intersects the line y = x.
It follows that
g(a) = f(a) −a > 0, and g(b) = f(b) −b < 0,

so that g(x) is a continuous function which is positive at a and negative at b, so by
the Intermediate Value Theorem (IVT), there exists c 2 (a, b) (open interval) with
g(c) = 0, i.e., f(c) = c, so c is a fixed point of f(x).  
Remark 1.2.8 The above is an example of an existence theorem, it says nothing
about how to find the fixed point, where it is or how many there are. It tells us that
if f(x) is a continuous function on an interval I with f(I)   I, then f(x) has a fixed
point in I. It is also true that if instead f(I)   I then f(x) has a fixed point in I as
we see now:
Theorem 1.2.9 Let f : I ! R (I = [a, b], a < b) be a continuous function with
f(I)   I, then f(x) has a fixed point in I.
Proof. As before, set g(x) = f(x) − x, then there exists c1 2 (a, b) with f(c1) < c1
(in fact f(c1) <a < c1). Also there is c2 2 (a, b) with f(c2) > c2.
Then g(c1) < 0 and g(c2) > 0 and since g(x) is a continuous function, it follows by
the Intermediate Value Theorem that there exists c 2 I, (c1 < c < c2 or c2 < c < c1),
with g(c) = 0; f(c) = c.  
Remark 1.2.10 It is often difficult to find fixed points explicitly:
Example 1.2.11 Set f(x) = cos x. If f(c) = c, then cos c = c. It is possible to find
an approximation to the fixed point c = .739085 . . . using (for example) Newton’s
method.

Exercises 1.2
1. Find all fixed points and eventual fixed points of the map f(x) = |x − 1|.
2. Let f : R ! R be such that for some n 2 Z+, the nth iterate of f has a unique
fixed point c (i.e., fn(c) = c and c is unique). Then show that c is a fixed point of f.
3. Use the Mathematica command NestList (see below) to find how the following
functions behave when the given point is iterated (comment on what appears to be
happening in each case):
(i) L1(x) = x(1 − x), with starting point x0 = .75.
(ii) L2(x) = 2x(1 − x), with starting point x0 = .1.
(iii) L3(x) = 3x(1 − x), with starting point x0 = .2.
(iv) L3.2(x) = 3.2x(1 − x), with starting point x0 = .95.
(v) f(x) = sin(x), with starting point x0 = 9.5.
(vi) g(x) = cos(x), with starting point x0 = −15.3.
4. Show that the logistic map Lμ(x) = μx(1 − x), x 2 [0, 1], for 0 < μ < 4, the fixed
point x = 0 has no eventual fixed points other than x = 1.
Show that for 1 < μ   2, the fixed point x = 1 − 1/μ only has x = 1/μ as an
eventual fixed point, but this is not the case for 2 < μ < 3.
To use Mathematica, first define your function, say f(x) = sin(x) in an input
cell using the syntax
f[x_]:= Sin[x].
Execute the cell by doing Shift and Enter simultaneously. The command
NestList[f, 9.5, 100]
when executed in a new input cell will give 100 iterations of sin(x) with starting
value x0 = 9.5. If you want to graph f(x) on the interval [a, b], use the command:
Plot[f[x], { x, a, b } ].

5. (a) Let f(x) = (1+x)−1. Find the fixed points of f and show that there are
no points c with f2(c) = c and f(c) 6= c (period 2-points). Note that f(−1) is not
defined, but points that get mapped to −1 belong to the interval [−2,−1) and are of
the form
 n = −
Fn+1
Fn
, n  1,
where {Fn}, n   0, is the Fibonacci sequence (see 1.1.4). Note that  n ! −r as
n!1, where −r is the negative fixed point of f (see [12] for more details).
(b) If x0 2 (0, 1], set xn =
Fn−1x0 + Fn
Fnx0 + Fn+1
. Use mathematical induction to show that
xn+1 = f(xn) =
Fnx0 + Fn+1
Fn+1x0 + Fn+2
.
Deduce that as n!1, xn ! 1/r, the positive fixed point of f.
1.3 Graphical Iteration
It is possible to follow the iterates of a function f(x) at a point x0 using graphical
iteration (sometimes called web diagrams).
We start at x0 on the x-axis and draw a line vertically to the function. We then
move horizontally to the line y = x, then vertically to the function and continue in this
way. Notice that in some examples the iterations converges to a fixed point. In others
it goes off to 1, whilst in others still, it oscillates between two points indefinitely.
Example 1.3.1 f(x) = x(1 − x). In this example, an examination of graphical
iteration seems to suggest that the orbits of any point in [0, 1] approach the fixed
point x = 0. This is an example of an attracting fixed point.

Example 1.3.2 Let f(x) = x/2 + 1. This is an affine transformation with a = 1/2,
b = 1. According to what we saw in Example 1.1.3, the iterates should converge to
b
1−a = 2 since |a| < 1. What is actually happening is that c = 2 is an attracting fixed
point of f(x) with the property that it attracts all members of R (said to be globally
attracting).
0 1 2 3 4
0
1
2
3
4
Example 1.3.3 We see from the graph of f(x) = sin x that c = 0 seems to be an
attracting fixed point. We shall show later that it is actually globally attracting.

We see two basic situations arising (together with some variations of these). In
particular, we notice that if the sequence xn = fn(x0) converges to some point c as
n!1, then c is a fixed point.
(i) Stable orbit, where the web diagram approaches a fixed point.
(ii) Unstable orbit, where the web diagram moves away from a fixed point.
0.0 0.2 0.4 0.6 0.8 1.0
0.0
0.2
0.4
0.6
0.8
1.0
0.5 1.0 1.5 2.0
0.5
1.0
1.5
2.0
Proposition 1.3.4 If y = f(x) is a continuous function and limn!1 fn(x0) = c,
then f(c) = c (i.e., if the orbit converges to a point c, then c is a fixed point of f) .
Proof. We see that
lim
n!1
fn(x0) = c ) f( lim
n!1
fn(x0)) = f(c),

and since f is continuous,
lim
n!1
fn+1(x0) = f(c).
But clearly c = limn!1 fn+1(x0), so that c = f(c) by the uniqueness of limit.  
1.4 Attractors and Repellers
Suppose that f : X ! X is a function (where we assume that X is some subset of
R such as an interval or possibly R itself). Let c 2 X be a fixed point of f. We define
the notions on attracting fixed point and repelling fixed that we looked at graphically
in the last section.
Definition 1.4.1 Let f : X ! X with f(c) = c.
(i) c is a stable fixed point if for all   > 0 there exists   > 0 such that if x 2 X and
|x − c| <  , then |fn(x) − c| <   for all n 2 Z+.
If this does not hold, c will be called unstable.
(ii) c is said to be attracting if there is a real number   > 0 such that
|x − c| <   ) lim
n!1
fn(x) = c.
(iii) c is asymptotically stable if it is both stable and attracting.
Remark 1.4.2 We will show that a fixed point c of f : X ! X (suitably differentiable)
with the property |f0(c)| < 1 is an asymptotically stable fixed point. If c
is an unstable fixed point, we can find   > 0 and x arbitrarily close to c such that
some iterate of x, say fn(x), will be greater than distance   from c. This is the case
when |f0(c)| > 1 as the next theorem shows. We call c a repeller (c is a repelling
fixed point), since the iterates move away from the fixed point (c is unstable). We will
also see that a fixed point can be stable without being attracting, and that it can be
attracting without being stable.
Definition 1.4.3 A fixed point c of f is hyperbolic if |f0(c)| 6= 1. If |f0(c)| = 1 it is
non-hyperbolic.
Thus if c is a non-hyperbolic fixed point, then f0(c) = 1 or f0(c) = −1, so the graph
of f(x) either meets the line y = x tangentially, or at 90 :

f(x) = −2x3 + 2x2 + x has both types of non-hyperbolic fixed points.
Hyperbolic fixed points have the following properties:
Theorem 1.4.4 Let f : X ! X be a differentiable function with continuous first
derivative (we say that f is of class C1).
(i) If a is a fixed point for f(x) with |f0(a)| < 1, then a is asymptotically stable.
The iterates of points close to a, converge to a geometrically (i.e., there is a constant
0 <   < 1 for which |fn(x) − a| <  n|x − a| for all n 2 Z+ and for all x 2 X
sufficiently close to a).
(ii) If a is a fixed point for f(x) for which |f0(a)| > 1, then a is a repelling fixed point
for f.
Proof. (i) We may assume that X is an open interval with a 2 X. Suppose |f0(a)| <
  < 1 for some   > 0, then using the continuity of f0(x), there exists an open interval
I with |f0(x)| <   < 1 for all x 2 I.
By the Mean Value Theorem, if x 2 I there exists c 2 I, lying between x and a,
satisfying
f0(c) =
f(x) − f(a)
x − a
,
so that
|f(x) − a| = |f0(c)||x − a| <  |x − a|,
(i.e., f(x) is closer to a than x was).
Repeating this argument with f(x) replacing x gives
|f2(x) − a| <  2|x − a|, . . . , |fn(x) − a| <  n|x − a|.

Since  n ! 0 as n!1, it follows that fn(x) ! a as n!1.
The proof of (ii) is similar.  
Example 1.4.5 Denote the Logistic map by Lμ(x) = μx(1 − x) = μx − μx2. We
are interested in this map for 0   x   1 and 0 < μ   4, since for these parameter
values μ, Lμ is a function which maps the interval [0, 1] into the interval [0, 1], so it
necessarily has a fixed point and iterates stay inside [0, 1] (if μ > 4, Lμ(x) > 1 for
some values of x in [0, 1] and further iterates will go to −1).
Solving Lμ(x) = x gives x = 0 or x = 1− 1/μ. If 0 < μ   1, then 1 − 1/μ   0, so
c = 0 is the only fixed point in [0, 1] in this case.
For 0 < μ   1, L0
μ(x) = μ − 2μx, so L0
μ(0) = μ and 0 is an attracting fixed point
for 0 < μ < 1.
0.0 0.2 0.4 0.6 0.8 1.0
0.2
0.4
0.6
0.8
1.0
If μ > 1, then 0 and 1 − 1/μ are both fixed points in [0, 1], but now 0 is repelling.

Also,
L0
μ(1 − 1/μ) = μ − 2μ(1 − 1/μ) = 2 − μ,
so that
|L0
μ(1 − 1/μ)| = |2 − μ| <1 iff 1< μ < 3.
In this case c = 1− 1/μ is an attracting fixed point, and repelling for μ > 3. Note
that L0
1(0) = 1 and L0
3(2/3) = 1 so we have non-hyperbolic fixed points.
We will examine in Chapter 2 what happens in the cases where μ = 1 and μ = 3.
Example 1.4.6 Let f(x) = 1/x. The two fixed points x = ±1 are hyperbolic since
f0(x) = −1/x2, so f0(±1) = −1. However, they are stable but not attracting since
f2(x) = f(1/x) = x, and we see points close to ±1 neither move closer nor further
apart.

Example 1.4.7 Consider instead the function fa(x) =
 
−2x if x < a
0 if x   a
, where
a is any positive real number. Then c = 0 is an unstable (repelling) fixed point of
fa which is attracting i.e., limn!1 fn(x) = 0 for all x 2 R. A fixed point having
this latter property is said to be globally attracting. Note that the function fa is not
continuous at x = a. It has been shown by Sedaghat [58] that a continuous mapping
of the real line cannot have an unstable fixed point that is globally attracting.
-3 -2 -1 0 1 2 3
-3
-2
-1
0
1
2
3
The map f1 (in bold) where x = 0 is an attracting but unstable fixed point.
Example 1.4.8 Newton’s Method Revisited.
Suppose that f : I ! I is a function whose zero c is to be approximated using
Newton’s method. The Newton function is Nf (x) = x − f(x)/f0(x), where we are
assuming f0(c) 6= 0. Notice that since f(c) = 0, Nf (c) = c, i.e., c is a fixed point of
Nf . Consider N0
f (c):
N0
f (x) = 1 −
f0(x)f0(x) − f(x)f00(x)
[f0(x)]2 =
f(x)f00(x)
[f0(x)]2 ,
so that
N0
f (c) =
f(c)f00(c)
[f0(c)]2 = 0,
since f(c) = 0.

It follows that |N0
f (c)| = 0 < 1, so that c is an attracting fixed point for Nf , and
in particular
lim
n!1
Nn
f (x0) = c,
provided x0, the first approximation to c, is sufficiently close to c.
Definition 1.4.9 A fixed point c for f(x) is called a super-attracting fixed-point if
f0(c) = 0. This gives a very fast convergence to the fixed-point for points nearby.
Remark 1.4.10 Suppose that f0(c) = 0, then Nf (x) is not defined at x = c, since the
quotient f(x)/f0(x) is not defined there. Suppose we can write f(x) = (x − c)kh(x)
where h(c) 6= 0 and k 2 Z+ (for example if f(x) is a polynomial with a multiple root),
then we have
f(x)
f0(x)
=
(x − c)kh(x)
k(x − c)k−1h(x) + (x − c)kh0(x)
=
(x − c)h(x)
kh(x) + (x − c)h0(x)
= 0,
when x = c, so that Nf (x) = c has a removable discontinuity at x = c (removed by
setting Nf (c) = c). Then we can find N0
f (x) and show that N0
f (x) has a removable
discontinuity at x = c, which can again be removed by setting N0
f (c) = (k − 1)/k, so
that |N0
f (c)| < 1 (see the exercises).
We summarize the above with a theorem:
Theorem 1.4.11 For a differentiable function f : I ! I, the zero c of f(x) is a
super-attracting fixed point of the Newton function Nf , if and only if f0(c) 6= 0.
Example 1.4.12 Suppose f(x) = x3 − 1, then f(1) = 0 and
Nf (x) = x − f(x)/f0(x) = x −
 
x3 − 1
3x2
 
=
2x
3
+
1
3x2 ,
so Nf (1) = 1 and N0
f (x) = 2
3
− 2
3x3 , so that N0
f (1) = 0. We see from graphical iteration
(next section), very fast convergence to the fixed point of Nf .

1. Find the fixed points and determine their stability for the function f(x) = 4 − 3
x .
2. For the family of quadratic maps Qc(x) = x2 + c, x 2 [0, 1], use the Mathematica
program webPlot below to give graphical iteration for the values shown (use 20
iterations):
(i) c = 1/2, starting point x0 = 1, (ii) c = 1/4, starting point x0 = .1, (iii) c = 1/8,
starting point x0 = .7.
Mathematica Program: Use as in:
webPlot[Cos[x], {x, 0, Pi}, {.1, 20}, AspectRatio->Automatic]
webPlot[g_, {x_, xmin_, xmax_}, {a1_, n_}, opts_] :=
Module[{seq, r, pts, web, graph}, r[t_] := N[g /. x -> t];
seq = NestList[r, a1, n]; pts = Flatten[
Table[{{seq[[i]],seq[[i + 1]]},{seq[[i + 1]],seq[[i + 1]]}},{i, 1, n}], 1];
web = Graphics[{Hue[0], Line[PrependTo[pts, {seq[[1]], 0}]]}];
graph = Plot[{x, r[x]}, {x, xmin, xmax},
DisplayFunction -> Identity]; Print["last iterate =", Last[seq]];
Show[web, graph, opts, Frame -> True,

PlotRange -> {{xmin, xmax}, {xmin, xmax}}]]
3. Let Sμ(x) = μ sin(x), 0   μ   2 , 0   μ     and Cμ(x) = μ cos(x), −    x    
and 0   μ    .
(a) Show that Sμ has a super-attracting fixed point at x =  /2, when μ =  /2.
(b) Find the corresponding values for Cμ having a super-attracting fixed point.
4. Show that the map f(x) =
2
x+ 1
has no periodic points of period n > 1 (Hint:
Use the closed formula in Exercise 1.1.6).
5. Show that if f(x) = x+1/x and x > 0, then Rn(x)!1as n!1. (Hint: R has
no fixed points in R and it is continuous with 0 < R(x) < R2(x) < · · · , so suppose
the limit exists and obtain a contradiction).
6. Let Nf be the Newton function of the map f(x) = x2 + 1. Clearly there are no
fixed points of the the Newton function as there are no zeros of f. Show that there
are points c where N2
f (c) = c (called period 2-points of Nf ).
7. (a) Suppose that f(c) = f0(c) = 0 and f00(c) 6= 0. If f00(x) is continuous at x = c,
show that the Newton function Nf (x) has a removable discontinuity at x = c (hint:
apply L’Hopital’s rule to Nf at x = c).
(b) If in addition, f000(x) is continuous at x = c with f000(c) 6= 0, show that N0
f (c) = 1/2,
so that x = c is not a super-attracting fixed point in this case.
(c) Check the above for the function f(x) = x3 − x2 with c = 0.
8. Continue the argument in Remark 1.4.11 (generalizing the last exercise), to show
that the derivative of the Newton function Nf has a removable discontinuity at x = c,
which can be removed by setting N0
f (c) = (k − 1)/k.

9. Let f be a twice differentiable function with f(c) = 0. Show that if we find the
Newton function of g(x) = f(x)/f0(x), then x = c will be a super-attracting fixed
point for Ng, even if f0(c) = 0 (this is called Halley’s method).
1.5 Non-hyperbolic Fixed Points.
Example 1.5.1We have seen that it f : X ! X (where usually X is some subinterval
of R) and a 2 X with f(a) = a, |f0(a)| < 1, with suitable differentiability conditions,
then a is a stable fixed point for f. This term is used because points close to a will
approach a under iteration. If we look at a function like f(x) = sin x, we see that
c = 0 is a fixed point but in this case |f0(0)| = 1, so it is a non-hyperbolic fixed point.
However, graphical iteration suggests that the basin of attraction of f is all of R, so
c = 0 is a stable fixed point. Before considering non-hyperbolic fixed points in more
detail, let us prove the last statement analytically:
Proposition 1.5.2 The fixed point c = 0 of f(x) = sin x is globally attracting .
Proof. First notice that c = 0 is the only fixed point of sin x. This is clear for if
sin x = x then we cannot have |x| > 1 since | sin x|   1 for all x. If 0 < x   1, the
Mean Value Theorem implies there exists c 2 (0, x) with
f0(c) =
f(x) − f(0)
x − 0
=
sin x
x
,
so that
sin x = x cosc < x,
since | cos c| < 1 for c 2 (0, 1). Similarly if −1   x < 0.
To show that f(x) is globally attracting, let x 2 R. We may assume that −1  
x   1, since this will be the case after the first iteration.
Suppose that 0 < x   1, then we note that 0 < f0(x) < 1 on this interval. By the
Mean Value Theorem, there exists c 2 (0, x) with
f0(c) =
f(x) − f(0)
x − 0
, or 0 < f(x) = f0(c)x < x.
Continuing in this way, we see that
0 < fn(x) < fn−1(x) < .. . < f(x) < x,
so we have a decreasing sequence xn = fn(x) bounded below by 0. It follows that this
sequence converges. But Proposition 1.4.5 implies that if this sequence converges, it

must converge to a fixed point. c = 0 being the only fixed point gives fn(x) ! 0 as
n!1. A similar argument can be used for −1   x < 0.  
Example 1.5.3 It is possible for the fixed point to be unstable, but to have a onesided
stability (to be semi-stable). For example, consider f(x) = x2 + 1/4 which has
the single (non-hyperbolic) fixed point c = 1/2. This fixed point is stable from the
left, but unstable on the right.
0.0 0.5 1.0 1.5 2.0
0.0
0.5
1.0
1.5
2.0
In the following we give some criteria for non-hyperbolic fixed points to be asymptotically
stable/unstable etc. It also gives a criteria for semi-stability.
Theorem 1.5.4 Let c be a non-hyperbolic fixed point of f(x) with f0(c) = 1. If f0(x),
f00(x) and f000(x) are continuous at x = c, then:
(i) if f00(c) 6= 0, the c is semi-stable,
(ii) if f00(c) = 0 and f000(c) > 0, then c is unstable,
(iii) If f00(c) = 0 and f000(c) < 0, then c is asymptotically stable.
Proof. (i) If f0(c) = 1 then f(x) is tangential to y = x at x = c. Suppose that
f00(c) > 0, then f(x) is concave up at x = c and the picture must look like the
following:

We see this gives stability on the left and instability on the right.
More formally, since the derivatives are continuous, and f00(c) > 0, this will be
true in some small interval (c −  , c +  ) surrounding c. In particular, the derivative
function f0(x) must be increasing on that interval, so that since f0(c) = 1 then
f0(x) < 1 for all x 2 (c −  , c), and f0(x) > 1 for all x 2 (c, c +  ),
for some   > 0. Also, from the continuity of f0(x), we can assume that f0(x) > 0 in
this interval.
Now by the Mean Value Theorem applied to the interval [x, c]   (c −  , c], there
exists q 2 (x, c) with
f0(q) =
f(x) − f(c)
x − c
,
Now since 0 < f0(q) < 1 and c > x, we have
0 <
f(x) − f(c)
x − c
< 1,
or
x < f(x) < c.
Repeating this argument, we see that the sequence fn(x) is increasing and bounded
above by c, so must converge to a fixed point. There can be no other fixed point (say
d 6= c), in this interval as the Mean Value Theorem would give f0(q1) = 1 for some
q1 2 (x, c), a contradiction. Consequently we see that c is stable on the left.
On the other hand, if [c, x]   (c, c+ ), then applying the Mean Value Theorem as
above gives
f0(q) =
f(x) − f(c)
x − c
> 1, so f(x) >x > c,

since x − c > 0. This tells us that the point moves away from c under iteration, so
the fixed point is unstable on the right. Similar considerations can be used when
f00(c) < 0 and the graph is concave down at x = c.
(ii) is similar to (iii), so is omitted.
(iii) In this case f000(c) < 0, f00(c) = 0 and f0(c) = 1. We will show that we have a
point of inflection at x = c as in the following picture:
0.5 1.0 1.5 2.0
0.5
1.0
1.5
2.0
By the second derivative test, f0(x) has a local maximum at x = c (the continuous
function f0(x) is concave down). It follows that
f0(x) < 1 for all x 2 (c −  , c +  ), x 6= c,
for some   > 0 (f00(x) > 0 for x 2 (c− , c), so f0(x) is increasing there, and f00(x) < 0
for x 2 (c, c +  ), so f0(x) is decreasing there). In particular f0(x) 6= 1 if x 6= c.
Now use an argument similar to that of (i) above to deduce the result.  
Example 1.5.5 Returning to the function f(x) = sinx, we see that f0(0) = 1,
f00(0) = 0 and f000(0) = −1, so the conditions of Theorem 1.5.4 (iii) are satisfied and
we conclude that x = 0 is an asymptotically stable fixed point.
If f(x) = tan x, then f0(0) = 1 and we can check that Theorem 1.5.4 (ii) holds
so that the fixed point x = 0 is unstable. For f(x) = x2 + 1/4, with f(1/2) = 1/2,
f0(1/2) = 1, we can apply Theorem 1.5.4 (i).

How do we treat the case where f0(c) = −1 at the fixed point? We use:
Definition 1.5.6 The Schwarzian derivative Sf(x) of f(x) is the function
Sf(x) =
f000(x)
f0(x)
−
3
2
 
f00(x)
f0(x)
 2
,
so that
Sf(x) = −f000(x) −
3
2
[f00(x)]2, when f0(x) = −1.
Theorem 1.5.7 Suppose that c is a fixed point for f(x) and f0(c) = −1. If f0(x),
f00(x) and f000(x) are continuous at x = c then:
(i) if Sf(c) < 0, then x = c is an asymptotically stable fixed point,
(ii) if Sf(c) > 0, then x = c is an unstable fixed point.
Proof. (i) Set g(x) = f2(x), then g(c) = c and we see that if c is asymptotically
stable with respect to g, then it is asymptotically stable with respect to f. Now
g0(x) =
d
dx
(f(f(x))) = f0(f(x)) · f0(x),
so that g0(c) = f0(c) · f0(c) = (−1)(−1) = 1.
The idea is to apply Theorem 1.5.4 (iii) to the function g(x). Now
g00(x) = f0(f(x)) · f00(x) + f00(f(x)) · [f0(x)]2,
thus
g00(c) = f0(c)f00(c) + f00(c)[f0(c)]2 = 0, since f0(c) = −1.
Also
g000(x) = f00(f(x))·f0(x)·f00(x)+f0(f(x))·f000(x)+f000(f(x))[f0(x)]3+f00(f(x))·2f0(x)f00(x).
Therefore
g000(c) = [f00(c)]2(−1) − f000(c) − f000(c) + 2f00(c)(−1)f00(c)
= −2f000(c) − 3[f00(c)]2
= 2Sf(c) < 0,
and the result follows from Theorem 1.5.4 (iii).
(ii) This now follows from Theorem 1.5.4 (ii)


Remark 1.5.8 The above proof shows how the Schwarzian derivative arises from
differentiating g = f   f = f2. In the case where f0(c) = −1, it follows that
g00(c) = 0 and Sf(c) =
1
2
g000(c).
Example 1.5.9 For the logistic map Lμ(x) = μx(1 − x) we have
L0
μ(x) = μ − 2μx, L00
μ(x) = −2μ and L000
μ (x) = 0.
When μ = 1, x = 0 is the only fixed point and Theorem 1.5.4 (i) shows that x = 0 is
semi-stable (attracting on the right). However, we regard this as a stable fixed point
for Lμ defined on the interval [0, 1], since points to the left of 0 are not in the domain
of Lμ.
When μ = 3, c = 2/3 is fixed and L0
μ(2/3) = −1 giving a non-hyperbolic fixed
point. However, Sf(2/3) = 0 − 3
2 [6]2 < 0 (negative Schwarzian derivative), so by
Theorem 1.5.7 (i), x = 2/3 is asymptotically stable.
0.0 0.2 0.4 0.6 0.8 1.0
0.0
0.2
0.4
0.6
0.8
1.0
Exercises 1.5
1. Show that f(x) = −2x3+2x2+x has two non-hyperbolic fixed points and determine
their stability.

2. For the family of quadratic maps Qc(x) = x2 + c, x 2 R, use the Theorems of
Section 1.5 to determine the stability of the fixed points for all possible values of c.
Find any values of c so that Qc has a non-hyperbolic fixed point, and determine their
stability.
3. Find the fixed points of the following maps and use the appropriate theorems to
determine whether they are asymptotically stable, semi-stable or unstable:
(i) f(x) =
x3
2
+
x
2
, (ii) f(x) = arctan x, (iii) f(x) = x3 + x2 + x,
(iv) f(x) = x3 − x2 + x, (v) f(x) =
 
0.8x ; if x   1/2
0.8(1 − x) ;if x > 1/2
4. Let f(x) = ax2 + bx + c, a 6= 0, and p a fixed point of f. Prove the following:
(i) If f0(p) = 1, then p is semistable.
(ii) If f0(p) = −1, then p is asymptotically stable.
5. If f(x) =
ax + b
cx + d
, a, b, c, d 2 R is the linear fractional transformation, show that
its Schwarzian derivative is Sf(x) = 0 for all x in its domain.
6. If Sf(x) is the Schwarzian derivative of f(x), a C3 function and F(x) =
f00(x)
f0(x)
,
show that Sf(x) = F0(x) − (F(x))2/2.
7. If f : R ! R is defined by f(x) =
 
x sin(1/x) if x 6= 0
0 ifx = 0
, find the fixed points
of f and show that for x 6= 0 they are non-hyperbolic. Show that x = 0 is not an
isolated fixed point (i.e., there are other fixed points arbitrarily close to 0). Is x = 0
a stable, attracting or repelling fixed point?

8. ([50]) Let Nf be the Newton function of a four times continuously differentiable
function f. If f( ) = 0, show that N000
f ( ) = 2Sf( ), where Sf is the Schwarzian
derivative of f.
9. (a) Use the Intermediate Value Theorem to show that f(x) = cos(x) has a fixed
point c in the interval [0,  /2]. We can show experimentally that this fixed point is
approximately c = .739085 . . ., for example by iterating any x0 2 R.
(b) Show that the basin of attraction of c is all of R (Hint: You may assume that
x 2 [−1, 1] - why? Now use the Mean Value Theorem to show that |f(x)−c| <  |x−c|
for some 0 <   < 1).
(c) Does f(x) have any eventual fixed points?
(d) Can f(x) have any points p with f2(p) = p other than c?




\end{document}

